\providecommand{\Ran}{\mathrm{Ran}}
\providecommand{\Fact}{\mathsf{Fact}}
\providecommand{\Perf}{\mathrm{Perf}}
\providecommand{\CoHA}{\mathrm{CoHA}}
\providecommand{\Prim}{\mathrm{Prim}}
\providecommand{\sdim}{\mathrm{sdim}}
\providecommand{\Spch}{\mathrm{Sp}^{\mathrm{ch}}}
\providecommand{\PhiFA}{\Phi^{\mathrm{FA}}}
\providecommand{\hCS}{\mathrm{hCS}}
\providecommand{\gDelta}{\mathfrak g_{\Delta_5}}

\section{Compact holomorphic \(E_3\) source from the volumes: attack--heal report}

\subsection{Claim attacked}

The attacked claim is the strongest possible compact-source claim:
the existing volumes already construct a compact holomorphic
\(E_3\)-factorization algebra
\[
        \mathcal A^{E_3}_{K3\times E}
\]
whose chain-level \(K3\)-to-\(E\) specialization, oriented compact
Hall object, protected integration, and primitive Lie comparison
produce the Igusa denominator algebra \(\gDelta\).

This claim is not proved in the volumes inspected.  The volumes contain
a coherent conditional architecture and several local, formal, quotient,
or character-level shadows.  They do not contain the compact
\(K3\times E\) source object together with the wrapped elliptic-degree
Ran repair, strong orientation, protected integration, and
presentation-level primitive comparison.  The strongest correct output
is therefore a construction route with named missing lemmas, not an
existence theorem.

\subsection{Source anchors with local file paths/lines}

\begin{itemize}
\item \verb|/Users/raeez/igusa-cusp-form/main.tex:164--183|:
the paper states the expected compact construction as a supplied
charge-completed holomorphic \(E_3\) factorization algebra followed by a
chain-level \(K3\)-to-\(E\) specialization; it also records the
projection-finite versus wrapped elliptic-degree obstruction.

\item \verb|/Users/raeez/igusa-cusp-form/main.tex:185--206|:
primitive recognition is explicitly conditional on a compact Igusa
datum, not a consequence of the scalar product.

\item \verb|/Users/raeez/igusa-cusp-form/main.tex:326--397|:
\(\Phi^{FA}_3(X)\) is a placeholder for an unconstructed chain-level
quantum object; the curve-level chiral shadow is only an expected
specialization.

\item \verb|/Users/raeez/igusa-cusp-form/main.tex:1908--1957|:
the holomorphic \(E_3\) quantum upgrade starts from a supplied
\(\mathcal A_X^{E_3}\); the \(\hCS\)/BV model is only a candidate after
descent, QME, anomaly, and completion checks.

\item \verb|/Users/raeez/igusa-cusp-form/main.tex:1981--2012|:
the chiral specialization is not invertible; the automorphic product is
a global trace and does not reconstruct local factorization maps.

\item \verb|/Users/raeez/igusa-cusp-form/main.tex:2014--2087|:
positive elliptic degree forces the support projection to be all of
\(E\), so naive support-local factorization on \(\Ran(E)\) misses the
Igusa \(s\)-direction.

\item \verb|/Users/raeez/igusa-cusp-form/main.tex:2179--2263|:
the compact Igusa realization datum lists the needed \(E_3\) source,
specialization, oriented Hall object, chiral coalgebra, cobar
comparison, and primitive recognition certificate.

\item \verb|/Users/raeez/igusa-cusp-form/main.tex:2265--2361|:
the primitive recognition theorem is conditional and does not construct
the compact object, bracket constants, parity lifts, simple primitive
representatives, or BKM relations.

\item \verb|/Users/raeez/igusa-cusp-form/main.tex:2408--2452|:
the determinant and Borcherds product determine signed exponents, not
Hall constants, parity dimensions, primitive representatives, relations,
or the Hopf radical quotient.

\item \verb|/Users/raeez/igusa-cusp-form/agent_material/12_compact_hall_ranE_attack_heal.tex:24--56|:
the OP/DT-to-Hall implication is attacked and repaired as a conditional
statement requiring an oriented critical Hall atlas, finite-first HN
completion, and Ran/wrapped locality theorem.

\item \verb|/Users/raeez/igusa-cusp-form/agent_material/12_compact_hall_ranE_attack_heal.tex:99--163|:
scalar stable-pair or Hilbert-scheme spaces are not an extension-closed
compact Hall stack, and the reduced \(E\)-quotient forgets relative
position data needed for chiral collisions.

\item \verb|/Users/raeez/igusa-cusp-form/agent_material/13_ranE_obstruction_attack_repair.tex:18--35|:
the positive elliptic-degree support lemma gives the precise wrapped
sector obstruction.

\item \verb|/Users/raeez/igusa-cusp-form/agent_material/13_ranE_obstruction_attack_repair.tex:58--85|:
the available repairs are either projection-finite local factorization
plus global Fock/Hecke wrapped modes, or a hybrid local/global wrapped
factorization formalism.

\item \verb|/Users/raeez/igusa-cusp-form/agent_material/12_primitive_comparison_attack_heal.tex:333--404|:
primitive comparison needs full presentation-level Hall and chiral
bar-cobar data; signed denominator multiplicities are insufficient.

\item \verb|/Users/raeez/igusa-cusp-form/agent_material/12_protected_integration_orientation_attack_heal.tex:232--309|:
orientation matching requires a determinant-line square root, equivariant
descent, protected integration, primitive projection, and the orientation
character \(\epsilon_o=\nu_{\Delta_5}\); constants do not supply it.

\item \verb|/Users/raeez/calabi-yau-quantum-groups/FRONTIER.md:27--31|:
Vol. III records the two-stage architecture
\(\Phi_d=\Spch\circ\PhiFA_d\), which is the correct schematic source for
the Igusa compact route.

\item \verb|/Users/raeez/calabi-yau-quantum-groups/FRONTIER.md:79--88| and
\verb|/Users/raeez/calabi-yau-quantum-groups/FRONTIER.md:197--211|:
the bracket-level BPS--BKM comparison and the Hall--Drinfeld
double/BKM-side object remain frontier problems.

\item \verb|/Users/raeez/calabi-yau-quantum-groups/chapters/theory/cy_to_chiral.tex:4--23|:
the formalism is explicitly two-stage: first an \(E_d\)-holomorphic
factorization algebra, then a chiral specialization.

\item \verb|/Users/raeez/calabi-yau-quantum-groups/chapters/theory/cy_to_chiral.tex:49--70|:
\(K3\times E\) has favorable formality input, but this is only a
Stage-1 formality witness, not the compact Hall source.

\item \verb|/Users/raeez/calabi-yau-quantum-groups/chapters/theory/cy_to_chiral.tex:90--104| and
\verb|/Users/raeez/calabi-yau-quantum-groups/chapters/theory/cy_to_chiral.tex:128--165|:
the specialization theorem requires Beck--Chevalley, Fubini/envelope,
proper-support, Tor-independence, compact-support, completion, framing,
and anomaly witnesses.

\item \verb|/Users/raeez/calabi-yau-quantum-groups/chapters/theory/bps_positive_geometry_closure.tex:231--247|:
the compact chamber certificate lists the exact Hall inputs: CY3
category, charge lattice, stability, support, orientation, critical
atlas, motivic target, HN finite control, integration, and realization
maps.

\item \verb|/Users/raeez/calabi-yau-quantum-groups/chapters/theory/bps_positive_geometry_closure.tex:286--330|:
the Igusa boundary statement is a quotient theorem against an automorphic
radical, not a proof that the raw compact Hall object exists with zero
radical.

\item \verb|/Users/raeez/calabi-yau-quantum-groups/chapters/theory/bps_positive_geometry_closure.tex:334--385|:
the \(\hCS\)-to-Hall descent is obstructed by explicit QME, anomaly,
gauge-fixing, critical-atlas, stationary-phase, and vertex
quasi-isomorphism classes.

\item \verb|/Users/raeez/calabi-yau-quantum-groups/chapters/examples/k3e_bkm_chapter.tex:13977--14073|:
the six \(K3\times E\) routes have conditional convergence; the common
character/denominator evidence is not yet a full bialgebra or compact
source theorem.

\item \verb|/Users/raeez/calabi-yau-quantum-groups/chapters/examples/k3e_bkm_chapter.tex:14082--14136|:
the numerical \(-\Delta_5^{-2}\) character is theorem-level, while its
identification as a compact critical CoHA character and full
M-theoretic/bialgebraic parent is conditional or conjectural.

\item \verb|/Users/raeez/chiral-bar-cobar/CLAUDE.md:56--63| and
\verb|/Users/raeez/chiral-bar-cobar/compute/lib/theorem_universal_chiral_genus_extension_engine.py:735--747|:
Vol. I starts with a supplied chiral or factorization algebra.  Its
bar-cobar and genus-extension mechanisms are downstream, not a compact
\(E_3\) source construction.

\item \verb|/Users/raeez/chiral-bar-cobar-vol2/CLAUDE.md:331--385|:
Vol. II separates the two-colour chiral/topological Swiss-cheese
structure from a single-colour \(E_3\) algebra and repeats that
Stage 2 is specialization, not inversion.

\item \verb|/Users/raeez/chiral-bar-cobar-vol2/chapters/theory/factorization_swiss_cheese.tex:35--58|:
local operadic shadows do not determine the global Ran object, periods,
monodromy, or global collision pattern.

\item \verb|/Users/raeez/chiral-bar-cobar-vol2/appendices/cfg_side_by_side.tex:659--755|:
the CFG/holographic rows concern \(\mathbb R^3\times K3\times\mathbb C^2\)
and a chiral avatar, not the compact \(K3\times E\) Hall/Ran source with
wrapped elliptic sectors.

\item \verb|/Users/raeez/topological-strings/main.tex:202--260| and
\verb|/Users/raeez/topological-strings/main.tex:1114--1165|:
the constructed topological-string sector is local and Hamiltonian; the
cochain-level factorization-center lift is stated as a problem.

\item \verb|/Users/raeez/topological-strings/main.tex:2238--2255| and
\verb|/Users/raeez/topological-strings/main.tex:2318--2322|:
the imported Costello--Li flat open-closed BCOV theorem does not prove a
compact \(K3\times E\) Hall or BKM source.
\end{itemize}

\subsection{Cross-volume lessons}

\begin{itemize}
\item Vol. I gives the \(E_1\)-chiral bar-cobar shadow once a chiral
algebra or factorization algebra is supplied.  It cannot be used
backwards to reconstruct a compact holomorphic \(E_3\) source.

\item Vol. II gives the Swiss-cheese chiral/topological local avatar and
records the local-to-global warning.  The local operad, formal
completion, and CFG/holographic \(E_3\) rows are not the compact
\(K3\times E\) construction.

\item Vol. III gives the correct architectural template:
\[
       \Perf(K3\times E)
       \xrightarrow{\ \PhiFA_3\ }
       E_3\text{-}\mathrm{HolFA}(K3\times E)
       \xrightarrow{\ \Spch_{K3,E}\ }
       \Fact(\Ran(E)).
\]
It also records why this template is conditional: the Stage-1 object,
Stage-2 specialization, compact Hall package, automorphic radical, and
bialgebra comparison still require witnesses.

\item The \(K3\times E\) examples in Vol. III support the denominator
and character side.  They do not identify the raw compact Hall object
before quotienting, and they do not prove that the automorphic radical
vanishes.

\item The topological-string volume supplies local BCOV/BV and
bulk-boundary intuition.  Its own text marks the cochain-level
factorization-center lift and compact CY output as unproved, so it is
evidence for the shape of the quantum-field-theoretic source, not a
compact Igusa construction.

\item The current Igusa paper is already aligned with these lessons:
the Borcherds product, \(-\Delta_5^{-2}\), and the formal current
envelope are trace, character, or denominator data.  They are not the
microscopic compact factorization algebra.
\end{itemize}

\subsection{First-principles construction candidate}

Let \(X=S\times E\), with \(S\) a \(K3\) surface.  The strongest honest
route is the following conditional construction.

\begin{enumerate}
\item Fix a charge lattice \(\Gamma\) containing the Igusa charges
\((n,\ell,m)\), a compact sectorial stability condition on
\(\Perf(X)\), and a finite-first HN completion.  Separate
projection-finite sectors from wrapped sectors by the projection
\(p_E\colon X\to E\).

\item Construct or supply a charge-completed holomorphic
\(E_3\)-factorization algebra
\(\mathcal A_X^{E_3}=\PhiFA_3(\Perf(X))\).  For a BV/\(\hCS\) model this
requires a strict \(E_3\) formality/framing witness, compact charge
completion, QME solution, anomaly cancellation, descent, and
compatibility with the chosen HN tower.

\item Construct the chain-level specialization
\[
       \Spch_{K3,E}(\mathcal A_X^{E_3})
          \simeq (p_E)_*\mathcal A_X^{E_3}
          =: \mathcal F_X
          \in \Fact(\Ran(E))
\]
with Beck--Chevalley, Fubini/envelope, Tor-independence, proper-support,
compact-support, and completion witnesses.  This gives an honest chiral
object only on the projection-finite part unless wrapped sectors are
added.

\item Repair the positive elliptic-degree obstruction by one of two
extra structures.  Route A is a projection-finite local Ran object
\(\mathcal F_X^{\mathrm{pf}}\) together with a global Fock/Hecke module
for classes with \(m>0\).  Route B is a hybrid wrapped factorization
formalism whose collision operations have both finite Ran legs and
global elliptic legs.  Either route must include associativity,
extension-stack functoriality, and compatibility with Hall
multiplication.

\item Build an oriented compact critical Hall object:
finite-type sectorial moduli stacks, extension correspondences,
critical/symplectic charts, vanishing cycles, strong orientation lines,
Thom--Sebastiani compatibility, \(E\)-equivariant descent, HN
wall-crossing, and protected integration.  The orientation character
must be identified with \(\nu_{\Delta_5}\), not chosen by normalization.

\item Compare Hall and chiral structures.  One needs a conilpotent
chiral coalgebra \(C_X\simeq \overline B^{\mathrm{ch}}_E(\mathcal F_X)\)
and a chiral cobar quasi-isomorphism to the formal denominator current
object \(\mathrm{Den}_{\Delta_5,E}\).

\item Prove the presentation-level primitive recognition certificate:
parity lifts of all root spaces, simple primitive representatives,
Borcherds--Kac--Moody relations, Hopf pairing, generation, and radical
quotient.  Only after this step may one conclude
\[
       H^0\Prim_{\mathrm{prot}}(\mathcal F_X)\cong \gDelta .
\]
\end{enumerate}

This is a construction candidate, not a theorem.  Every arrow above is
load-bearing; omitting any one of them collapses the argument to a
character or denominator shadow.

\subsection{Exact missing lemmas}

\begin{enumerate}
\item A compact Stage-1 lemma constructing
\(\mathcal A_{K3\times E}^{E_3}\) from \(\Perf(K3\times E)\), with
strict \(E_3\) formality, framing, charge completion, and descent
witnesses.

\item A compact BV/\(\hCS\) QME and anomaly lemma for \(K3\times E\),
not only for flat or noncompact local models.

\item A finite-first HN lemma for the chosen compact charge sectors,
compatible with Hall extension correspondences and the completed
\(E_3\) algebra.

\item A chain-level specialization lemma proving
\(\Spch_{K3,E}(\mathcal A_X^{E_3})\simeq (p_E)_*\mathcal A_X^{E_3}\)
with Beck--Chevalley, Fubini, proper-support, Tor-independence, and
completion hypotheses.

\item A wrapped elliptic-degree locality lemma.  It must either build
the global Fock/Hecke mode coupled to projection-finite Ran collisions,
or define a hybrid wrapped factorization category and prove its
associativity.

\item An extension-closed compact Hall-heart lemma for the relevant
BPS/stable-object category on \(K3\times E\), including stability,
support, wall-crossing, and finite-type sectorial moduli.

\item A critical atlas and vanishing-cycle lemma giving the required
shifted symplectic/critical charts on the compact Hall stacks.

\item A strong orientation lemma constructing determinant-line square
roots compatible with extension, Thom--Sebastiani, \(E\)-descent, and
HN wall crossing.

\item An orientation-monodromy lemma proving
\(\epsilon_o=\nu_{\Delta_5}\) on the type-II Weyl group from determinant
line transport.

\item A protected integration lemma whose primitive signed
superdimensions are the Jacobi coefficients \(f(nm,\ell)\), with
parity lifts rather than only signed dimensions.

\item A Hall-to-BKM primitive lemma proving the simple primitive
representatives, bracket constants, Serre/Borcherds relations,
generation, Hopf pairing, and radical quotient.

\item A chiral Koszul comparison lemma proving
\[
       C_X\simeq \overline B^{\mathrm{ch}}_E(\mathcal F_X),
       \qquad
       \Omega^{\mathrm{ch}}_E C_X\simeq \mathrm{Den}_{\Delta_5,E},
\]
with conilpotence and completed-current control.
\end{enumerate}

\subsection{Falsified shortcuts}

\begin{itemize}
\item The scalar OP/DT identity does not construct a compact Hall
factorization algebra.  Stable-pair or Hilbert-scheme spaces give
enumerative invariants, not an extension-closed Hall stack with
collision maps.

\item The Borcherds product and denominator identity do not determine
the compact \(E_3\) source.  They determine signed root
supermultiplicities; a zero bracket on the same signed dimensions has
the same product and the wrong Lie algebra.

\item The formal current envelope \(\mathrm{Den}_{\Delta_5,E}\) is not a
microscopic source.  It is a target or comparison object after the
compact factorization/Hall data have been built.

\item The Vol. II CFG/holographic \(E_3\) rows on
\(\mathbb R^3\times K3\times\mathbb C^2\) do not give the compact
\(K3\times E\) source.  They lack the compact elliptic wrapping,
sectorial Hall stacks, and orientation package required here.

\item A local Swiss-cheese or formal operadic shadow does not determine
the global Ran object.  Global periods, monodromy, and wrapped elliptic
degree are invisible to the local operad.

\item \(K3\times E\) formality and Atiyah-class vanishing do not by
themselves solve compact anomaly cancellation, charge completion,
specialization, or Hall comparison.

\item The six \(K3\times E\) routes and the four \(\kappa\)-values in
Vol. III are not six applications of one functor producing the same
object.  Their convergence is conditional on the missing Hall/BKM
package.

\item A quasi-NCCR or boundary character theorem is not a global NCCR
and not the raw compact Hall theorem.  It can fix the automorphic
normalization while quotienting out classes invisible to the
denominator.

\item The local BCOV/topological-string sector does not prove compact
\(K3\times E\) Hall or BKM realization.  The topological-string volume
itself marks the relevant cochain-level center and compact CY outputs as
open.

\item The constants \(64\) and \(-4096\) do not determine orientation.
Orientation is determinant-line transport with monodromy; it is not a
normalization choice.
\end{itemize}

\subsection{Recommended repair theorem/open problem language}

\paragraph{Conditional theorem.}
Assume the following data for \(X=K3\times E\):
a charge-completed holomorphic \(E_3\)-factorization algebra
\(\mathcal A_X^{E_3}\) with QME, anomaly, descent, framing, and compact
completion witnesses; a chain-level specialization
\(\Spch_{K3,E}(\mathcal A_X^{E_3})=\mathcal F_X\) satisfying the
Beck--Chevalley/Fubini/proper-support/Tor/completion hypotheses; a
projection-finite plus wrapped-sector Ran formalism; an oriented
HN-completed compact critical Hall object with protected integration and
\(\epsilon_o=\nu_{\Delta_5}\); a conilpotent chiral bar-cobar comparison
to \(\mathrm{Den}_{\Delta_5,E}\); and a full primitive recognition
certificate.  Then
\[
       H^0\Prim_{\mathrm{prot}}(\mathcal F_X)\cong \gDelta
\]
as a Borcherds--Kac--Moody Lie superalgebra, and its Euler/supertrace
denominator is the Igusa \(\Delta_5\) product.

\paragraph{Open problem.}
Construct the above compact Igusa realization datum.  Equivalently,
prove the missing lemmas listed in this report: compact holomorphic
\(E_3\) source, compact BV anomaly/QME control, chain-level
\(K3\)-to-\(E\) specialization, wrapped elliptic Ran repair, oriented
compact Hall stacks, protected integration, orientation monodromy,
primitive presentation, and chiral bar-cobar comparison.  Until those
data are constructed, the correct claim is conditional recognition from
a supplied compact source, not existence of the source.

\subsection{Files changed}

\begin{itemize}
\item \verb|/Users/raeez/igusa-cusp-form/agent_material/14_compact_E3_source_from_volumes_attack_heal.tex|
\end{itemize}

\subsection{Checks run}

\begin{itemize}
\item Read the repository instructions and manuscript anchors:
\verb|CLAUDE.md|, \verb|AGENTS.md|, \verb|main.tex|, and \verb|proj.bib|.

\item Read the relevant prior attack-heal reports under
\verb|agent_material/12_*| and \verb|agent_material/13_*|, especially the
compact Hall/Ran(E), primitive comparison, protected integration, and
Ran(E) obstruction reports.

\item Searched and read the cross-volume anchors in
\verb|/Users/raeez/chiral-bar-cobar|,
\verb|/Users/raeez/chiral-bar-cobar-vol2|,
\verb|/Users/raeez/calabi-yau-quantum-groups|, and
\verb|/Users/raeez/topological-strings| listed above.

\item Checked the active worktree before writing.  Existing unrelated
modifications to \verb|appendices/boundary_compatibility_conditions.tex|,
\verb|main.tex|, and \verb|out/main.pdf| were left untouched.

\item No manuscript build was run for this report-only change; building
would not test a standalone agent report and could mutate shared output
files.
\end{itemize}
